\documentclass[10pt,a4paper]{amsart}
\usepackage[margin=19mm]{geometry}
\usepackage{amsmath,amssymb,mathtools}
\usepackage[scr=rsfs,cal=euler]{mathalfa}
\usepackage{xfrac}
\usepackage[unicode,hidelinks,bookmarks=false]{hyperref}
\newcommand{\ie}{i.e.\ }
\newcommand{\cf}{cf.\ }
\newcommand{\resp}{respectively\ }
\newcommand{\Subsection}{\subsection}
\providecommand{\nobreakdash}{\nobreak\hbox{-}}
\setlength{\emergencystretch}{2em}
\multlinegap0pt
\usepackage{bm}
\usepackage{bbm}
\usepackage{dsfont}
\usepackage{xspace}
\usepackage{cancel}
\usepackage{mathtools}
\usepackage{amsthm}
\makeatletter
\@ifundefined{lemma}{\newtheorem{lemma}{Lemma}}{}
\@ifundefined{proposition}{\newtheorem{proposition}{Proposition}}{}
\makeatother
\newcommand{\N}{\mathbb{N}}
\newcommand{\Z}{\mathbb{Z}}
\title[On variants of Chowla's conjecture]{On variants of Chowla's conjecture}
\author{Krishnarjun Krishnamoorthy}
\date{Source: arXiv:2501.10962v2 (CC BY 4.0)}
\begin{document}
\maketitle

\noindent\emph{Source and licence.} On variants of Chowla's conjecture. Version arXiv:2501.10962v2 (CC BY 4.0), distributed under the Creative Commons Attribution 4.0 International licence, as recorded in the arXiv licence metadata for this exact version and shown on the abstract page https://arxiv.org/abs/2501.10962v2. Full rights notice: licenses/arxiv-2501.10962-CC-BY-4.0.txt.\par\smallskip

\noindent\textit{Editorial scope.} Counted unit: the Proposition whose source label is \texttt{Proposition "Approximation"} in the article (source0.tex lines 327--336, source label \texttt{Proposition "Approximation"}), delivered together with the article's own complete original proof of it (source0.tex lines 338--402, one \texttt{proof} environment of 65 lines). No mathematics is written, repaired or re-derived: statement and proof are the article's own text line for line. Required context retained uncounted: Lemma "N\_p recursion" (lines 241--246); Lemma "Properties of AP" (lines 283--291); Proposition "Multiplicative AP" (lines 306--308). Every definition, display and label of those lines is retained unchanged. Omitted from this package and not redistributed: 1--240, 247--282, 292--305, 309--326, 337--337, 403--530. Nothing inside a retained excerpt is omitted or altered: every retained body is delivered exactly as the article's lines are written, so no omission receipt of any kind accompanies this package. Citations: the retained excerpts carry no citation command at all, so no bibliography entry of the article is transcribed and no reference is redistributed. Reference adaptation: none was needed and none was made; every reference inside the delivered unit resolves inside it, so no unresolved reference or citation is produced. Printed slips kept literal: no printed word, symbol, hypothesis or number of the counted statement or of its proof was altered. Layout and class: the article's own class is \texttt{amsart} with packages including amsmath, amsfonts, amssymb, tikz, amsthm, hyperref, inputenc, babel, xcolor, caption, epigraph, mathtools; the wrapper is the lane's pinned \texttt{amsart} editorial wrapper at 10pt on A4, which restates the theorem-like environments the retained text needs and adds the article's own macro definitions that the retained text needs (\texttt{\textbackslash N}, \texttt{\textbackslash Z}), with extra wrapper packages (bm, bbm, dsfont, xspace, cancel, mathtools). The delivered numbering is the unit's own. Assets: no retained span contains a figure environment, a graphics-inclusion command, a PDF-page inclusion, a table or a TikZ picture; no image file is copied into this package, the package declares no asset dependency and no image file is redistributed with it. Licence: version arXiv:2501.10962v2 of this article is distributed under the Creative Commons Attribution 4.0 International licence, as recorded in the arXiv licence metadata for this exact version and shown on the abstract page https://arxiv.org/abs/2501.10962v2. A full rights notice is delivered at licenses/arxiv-2501.10962-CC-BY-4.0.txt. Characters: every retained excerpt is pure ASCII, so the delivered engine, XeLaTeX with its default Latin Modern text font, reports no missing character.
	\begin{lemma}\label{Lemma "N_p recursion"}
		For any two subsets $ P_1, P_2 $ of the primes and two finite subsets $H_1,H_2$ of $\N_0$,
		\[
		\mathcal{N}_{P_1\triangle P_2}^{H_1\triangle H_2} = \mathcal{N}_{P_1}^{H_1} \triangle \mathcal{N}_{P_1}^{H_2} \triangle \mathcal{N}_{P_2}^{H_1} \triangle \mathcal{N}_{P_2}^{H_2}.
		\]
	\end{lemma}

	\begin{lemma}\label{Lemma "Properties of AP"}
		For any non-empty subset $P$ of the primes, the following statements are true.
		\begin{enumerate}
			\item	$\mathcal{A}_P$ is closed under intersection.
			\item 	$\mathcal{A}_P$ is closed under complements.
			\item 	Every element of $\mathcal{A}_P$ can be written as a finite disjoint union of sets of the form $\{C(r,n)\}$ for a fixed $n\in \mathcal{S}_P$.
			\item 	The natural density exists for every set $A\in \mathcal{A}_P$.
		\end{enumerate}
	\end{lemma}

	\begin{proposition}\label{Proposition "Multiplicative AP"}
		 Suppose that $P_1, P_2$ are non-empty mutually disjoint subsets of primes. Then $\delta(A\cap B) = \delta(A)\delta(B)$ for any $A\in \mathcal{A}_{P_1}$ and $B\in \mathcal{A}_{P_2}$.
	\end{proposition}

	 \begin{proposition}\label{Proposition "Approximation"}
	 	Given a finite subset $P$ of the primes and $H\subseteq \N_0$ and an integer $r$, there exists subsets $X_P^H(r), Y_P^H(r)$ in $\mathcal{A}_P$ such that 
	 	\[
	 	X_P^H(r)\subseteq \mathcal{N}_P^H\subseteq Y_P^H(r)
	 	\]
	 	and 
	 	\[
	 	\lim_{r\to\infty} \delta(X_P^H(r)) = \lim_{r\to\infty} \delta(Y_P^H(r)).
	 	\]
	 \end{proposition}

	\begin{proof}
		The proof is by double induction, first on $|H|$ and then on $|P|$. Let us first prove the statement for a fixed prime $p$ and when $H=\{h\}$. Then $\mathcal{N}_p^h$ contains precisely those integers $n$ such that an odd power of $p$ properly divides $n+h$. In particular,
		\begin{equation}\label{Equation "N_p^h union"}
			\mathcal{N}_p^h = \bigcup_{r=1}^{\infty} \left(C(h, p^{2r-1})\setminus C(h, p^{2r})\right).
		\end{equation}
		Since the above union is a disjoint union, we have the inequalities,
		\begin{equation}\label{Equation "X_p, Y_p 1"}
			\bigcup_{t=1}^{r} \left(C(h, p^{2t-1})\setminus C(h, p^{2t})\right) \subseteq \mathcal{N}_p^h \subseteq \bigcup_{t=1}^{r}\left( C(h, p^{2t-1})\setminus C(h, p^{2t})\right) \bigcup C(h, p^{2r-1})
		\end{equation}
		for any $r$. Call the sets on the left hand side and right hand side of \eqref{Equation "X_p, Y_p 1"} as $X_p^h(r)$ and $Y_p^h(r)$ respectively. We observe that $\delta(X_p^h(r))\leqslant \delta(Y_p^h(r))\leqslant \delta(X_p^h(r)) + p^{-2r+1}$. Letting $r\to\infty$, we see that $\delta(\mathcal{N}_p^h)$ exists and equals
		\begin{equation}\label{Equation "Temp 1"}
			\delta\left(\mathcal{N}_p^h\right) = \sum_{r=1}^{\infty} \left(\frac{1}{p^{2r-1}} - \frac{1}{p^{2r}}\right)= \frac{1}{p+1}.
		\end{equation}
		This completes the proof in this case.
		
		Suppose that $H= \{h_1,\ldots,h_d\}$ and consider $\mathcal{N}_p^H$. Let $\varphi : H\to \Z/p\Z$ denote the natural projection map. From Lemma \ref{Lemma "N_p recursion"},
		\begin{equation}\label{Equation "Temp"}
			\mathcal{N}_p^H =  \mathcal{N}_p^{\varphi^{-1}([0])}\triangle\ldots \triangle \mathcal{N}_p^{\varphi^{-1}([p-1])} = \mathcal{N}_p^{h_1}\triangle\ldots \triangle \mathcal{N}_p^{h_d}
		\end{equation}
		In particular, from \eqref{Equation "Temp 1"}, it follows that 
		\begin{equation}\label{Equation "eta_p evaluation"}
			\delta^+\left(\mathcal{N}_p^H\right)\leqslant \frac{d}{p+1}.
		\end{equation}
		But from \eqref{Equation "N_p^h union"}, we may easily deduce that $\mathcal{N}_p^{h_i}\cap \mathcal{N}_p^{h_j}\neq \phi$ only if $p| (h_i-h_j)$. Therefore we may rewrite the first equality in \eqref{Equation "Temp"} as the disjoint union
		\begin{equation}\label{Equation "Disjoint union"}
			\mathcal{N}_p^H = \bigcup_{i=0}^{p-1} \mathcal{N}_p^{\varphi^{-1}([i])}.
		\end{equation}
		If $\varphi$ is a non-constant function, then $\eta_p^H$ exists by induction on $|H|$ and we are done. In particular, we may suppose that $H\neq \{0,1\},\{1,2\}$ or $\{0,1,2\}$. 
		
		Suppose that $\varphi$ is a constant function, say $\varphi(H) = \{[i_1]\}$ for some $0\leqslant i_1\leqslant p-1$. If $n\not\equiv -i_1\mod p$, then $p\nmid (n+h)$ for any $h\in H$ and hence $\Lambda_p^H(n)=1$. Therefore, $n\in \mathcal{N}_p^H$ only if $n\equiv -i_1\mod p$. Moreover, in this situation,
		\[
		\Lambda_p^H(n) = \lambda_p(n+h_1)\ldots\lambda_p(n+h_d) = (-1)^d \lambda_p\left(\frac{n}{p} + \frac{h_1}{p}\right)\ldots \lambda_p\left(\frac{n}{p} + \frac{h_d}{p}\right).
		\]
		Let $H_1:=\frac{1}{p}(H-i_1)$. We observe that $\max\{H_1\} < \max\{H\}$. Rewriting the above,
		\begin{equation}\label{Equation "Calculation"}
			\mathcal{N}_p^H = \begin{cases}
				p \mathcal{N}_p^{H_1} - i_1&\mbox{if }2| d\\
				p \left(\mathcal{N}_p^{H_1}\right)^c - i_1&\mbox{otherwise.}
			\end{cases}
		\end{equation}
		In either case, it is enough to prove the proposition for $\mathcal{N}_p^{H_1}$. Now, if we consider the natural projection map $\varphi_1: H_1\to \Z/p\Z$, we may argue as above. If $\varphi_1$ is a non-constant map, we are done. Otherwise we obtain (if necessary, after translation and scaling as above) $H_2:= \frac{1}{p} (H_1-i_2)$, and the problem reduces to that of $\mathcal{N}_p^{H_2}$, with $\max\{H_2\} < \max\{H_1\}$. Proceeding forward this process ultimately terminates because $\max\{H_i\}$ is strictly decreasing. In particular, there is a stage $r$ where the projection map (say $\varphi_r : H_r\to \Z/p\Z$) is non-constant and we may then apply induction on $|H|$. This completes the proof when $|P|=1$. We note in passing that for large enough primes $p$, $\varphi$ is an injection and hence from \eqref{Equation "Temp 1"} and \eqref{Equation "Disjoint union"} we have $\eta_p^H = \frac{d}{p+1}.$
		
		Now fix $H$ and suppose that $P = \{p_1,p_2,\ldots, p_k\}$. We shall prove the proposition holds for $P$ supposing the same for $P':=\{p_1,p_2,\ldots, p_{k-1}\}$ and $\{p_k\}$. From Lemma \ref{Lemma "N_p recursion"}, 
		\[
		\mathcal{N}_P^H = \mathcal{N}_{P'}^H\triangle \mathcal{N}_{\{p_k\}}^H.
		\]
		For every $r$, choose $X_{P'}^H(r), Y_{P'}^H(r), X_{\{p_k\}}^H(r), Y_{\{p_k\}}^H(r)$ from induction. Then, we have
		\[
		\left(X_{P'}^H(r)\setminus Y_{\{p_k\}}^H(r)\right) \bigcup \left(X_{\{p_k\}}^H(r) \setminus Y_{P'}^H(r)\right) \subseteq \mathcal{N}_P^H \subseteq \left(Y_{P'}^H(r)\setminus X_{\{p_k\}}^H(r)\right) \bigcup \left(Y_{\{p_k\}}^H(r) \setminus X_{P'}^H(r)\right).
		\]
		From Lemma \ref{Lemma "Properties of AP"}, both the left hand side and the right hand side above belong to $\mathcal{A}_P$. In particular the natural densities exist for each of those sets. Considering (upper) natural densities (and observing that $X_{P'}^H(r)\subseteq Y_{P'}^H(r)$) we get,
		\begin{multline*}
			\delta\left(X_{P'}^H(r)\setminus Y_{\{p_k\}}^H(r)\right) +  \delta\left(X_{\{p_k\}}^H(r) \setminus Y_{P'}^H(r)\right)  \leqslant \delta^+(\mathcal{N}_P^H )\\
			\leqslant \delta\left(Y_{P'}^H(r)\setminus X_{\{p_k\}}^H(r)\right) + \delta \left(Y_{\{p_k\}}^H(r) \setminus X_{P'}^H(r)\right).
		\end{multline*}
		From Proposition \ref{Proposition "Multiplicative AP"}, we get
		\begin{multline*}
			\delta(X_{P'}(r))(1-\delta(Y_{\{p_k\}}(r))) + \delta(X_{\{p_k\}}(r))(1-\delta(Y_{P'}(r))) \leqslant \delta^+(\mathcal{N}_P) \\
			\leqslant \delta(Y_{P'}(r))(1-\delta(X_{\{p_k\}}(r))) + \delta(Y_{\{p_k\}}(r))(1-\delta(X_{P'}(r))).
		\end{multline*}
		Similarly, we may also consider lower natural densities. Letting $r\to\infty$, we prove $\delta(\mathcal{N}_P^H)$ exists and equals
		\begin{equation}\label{Equation "eta_P recursion"}
			\eta_{P}^H = \eta_{P'}^H(1-\eta_{p_k}^H) + \eta_{p_k}^H(1-\eta_{P'}^H).
		\end{equation}
	\end{proof}

\end{document}
